\documentclass[english, parskip=full]{scrartcl}
\usepackage[utf8]{inputenc}  % Kodierung der Datei
\usepackage[ngerman]{babel}  % Sprache auf Deutsch 
\usepackage{color}
\usepackage{graphicx, gensymb}
\usepackage{amsfonts, cancel, amssymb}
\usepackage{amsmath, amsthm, array, esint}
\usepackage{booktabs, multirow}  % schönere Tabellen
\usepackage{caption}  % Anpassung der Captions
\usepackage{scrlayer-scrpage}    % Manipulation des Seitenstils
\pagestyle{scrheadings}  % Stil für die Seite setzen
\automark{section}  % Abschnittsnamen als Seitenbeschriftung 
\setheadsepline{.5pt}    % Kopzeile durch Linie abtrennen
\usepackage[hidelinks]{hyperref}  % Links und weitere PDF-Features
\usepackage{typearea, tikz, pgffor, verbatim}
\usetikzlibrary{calc, arrows.meta,arrows, graphs, quotes, positioning}
\usepackage[cache=false, outputdir=build]{minted}
\usemintedstyle{vs}
\usepackage{placeins} % provides \FloatBarrier

\definecolor{bg}{rgb}{0.9,0.9,0.9}
\newcommand\numberthis{\addtocounter{equation}{1}\tag{\theequation}}
\usepackage{svg}
\usepackage{siunitx}
\usepackage{physics}
\usepackage{tensor}
\renewcommand{\bra}{}
\newcommand{\kom}{}
\newcommand{\brak}{}
\renewcommand{\braket}{}
\newcommand{\intx}{}
\newcommand{\intr}{}
\newcommand{\kla}{}
\newcommand{\klam}{}
\newcommand{\klammer}{}
\newcommand{\oper}{}
\newcommand{\tecxt}{}
\newcommand{\Bra}{}
\newcommand{\Ket}{}
\renewcommand{\i}{\text{i}}
\newcommand{\e}{\text{e}}
\newcommand*{\Fourier}{\textsc{Fourier}\ }
\newcommand*{\F}{\mathcal{F}}
\newcommand*{\sinc}{\text{sinc\,}}
\newcommand{\conv}{\mathop{\scalebox{3.5}{\raisebox{-0.38ex}{\makebox[0.5\width][c]{$\ast$}}}}\limits}

\newcommand*{\titel}{Bandstrukturen in verschiedenen Modellen}
\newcommand*{\autor}{Joachim Schwardt}

\hypersetup{pdfauthor={\autor}, pdftitle={\titel}}

%\titlehead{titlehead}
\subject{Theoretische Festkörperphysik}
\title{\titel}
\author{\autor}
\date{\today}

\begin{document}

%\begin{titlepage}
\maketitle  % Titel setzen
%\tableofcontents  % Inhaltsverzeichnis setzen
%\end{titlepage}

\section{Quadratisches Gitter}
%%%% BRILLOUIN ZONE CONSTRUTION !!!!! 
%%%% https://www.doitpoms.ac.uk/tlplib/brillouin_zones/zone_construction.php
Gegeben ist ein quadratisches Gitter mit der Gitterkonstanten $a$.
Für ein freies Elektron gilt dann
\begin{align}
\epsilon_n(\textbf{k}) &= \frac{\hbar^2}{2m} \kla\left( \textbf{k} - \textbf{G} \right)^2 = \frac{\hbar^2}{2m} \klam\left[ \kla\left( k_x - \frac{2\pi}{a} n_x \right)^2 + \kla\left( k_y - \frac{2\pi}{a}n_y \right)^2 \right].
\end{align}
%Die erste Brillouin-Zone besteht hier aus allen $\textbf{k}$ mit $|k_{x,y}| \le \frac{\pi}{a}$.
%Insbesondere gilt in der ``Ecke'' also
%\begin{align}
%\epsilon_n \kla\left( \frac{\pi}{a}, \frac{\pi}{a} \right) &= \frac{\hbar^2}{2m} \frac{\pi^2}{a^2} \klam\left[ \kla\left( 1 - 2n_x \right)^2 + \kla\left( 1 - 2n_y \right)^2 \right] =  \\
%&= 2 \frac{\hbar^2\pi^2}{2ma^2} \klam\left[ 1 + 2(n_x^2+n_y^2) - 2(n_x+n_y) \right]
%\end{align}
%und im Mittelpunkt einer Seitenlinie der ersten Brillouin-Zone gilt
%\begin{align}
%\epsilon_n \kla\left( \frac{\pi}{a}, 0 \right) &= \frac{\hbar^2\pi^2}{2ma^2} \klam\left[ (1 - 2n_x)^2 + (2n_y)^2 \right] = \\
%&= \frac{\hbar^2\pi^2}{2ma^2} \klam\left[ 1 - 4n_x + 4n_x^2 + 4n_y^2 \right].
%\end{align}
%Für das erste Band ist $n_x=n_y=0$, womit also $\epsilon_0\kla\left( \frac{\pi}{a},\frac{\pi}{a} \right) = 2 \epsilon_0 \kla\left( \frac{\pi}{a}, 0 \right)$.
Wir betrachten hier das ausgedehnte Zonenschema, es gibt also effektiv nur ein Band, was dafür allerdings nicht in $\textbf{k}$ begrenzt ist.
Die erste Brillouin-Zone besteht hier aus allen $\textbf{k}$ mit $|k_{x,y}| \le \frac{\pi}{a}$.
Insbesondere gilt in der ``Ecke'' also
\begin{align}
\epsilon \kla\left( \frac{\pi}{a}, \frac{\pi}{a} \right) &= \frac{\hbar^2}{2m} \klam\left[ \kla\left( \frac{\pi}{a} \right)^2 + \kla\left( \frac{\pi}{a} \right)^2 \right] = \frac{\hbar^2\pi^2}{ma^2}
\end{align}
und im Mittelpunkt einer Seitenlinie der ersten Brillouin-Zone gilt
\begin{align}
\epsilon \kla\left( \frac{\pi}{a}, 0 \right) &= \frac{\hbar^2}{2m} \left[ \kla\left( \frac{\pi}{a} \right)^2 + 0 \right] = \\
&= \frac{\hbar^2\pi^2}{2ma^2}.
\end{align}
Also gilt $\epsilon_0\kla\left( \frac{\pi}{a},\frac{\pi}{a} \right) = 2 \epsilon_0 \kla\left( \frac{\pi}{a}, 0 \right)$.
\def\bw{blue!50!white}
\def\ow{orange!80!white}
\begin{figure}[!ht]
\centering
\begin{tikzpicture}[scale=2]   
    \pgfmathsetmacro{\r}{0.03}
    \pgfmathsetmacro{\xmin}{-2}
    \pgfmathsetmacro{\xmax}{2}
    \pgfmathsetmacro{\ymin}{-2}
    \pgfmathsetmacro{\ymax}{2}

    % 1. BZ
    \filldraw[fill=blue!25!white, draw=black] (-0.5,-0.5) --++ (1,0) --++ (0,1) --++ (-1,0) -- cycle;
    \foreach \v in {-0.5,0.5}{
        \draw[\bw, thick] (\v,\ymin) -- (\v,\ymax);
        \draw[\bw, thick] (\xmin,\v) -- (\xmax,\v);
    }
    \foreach \x in {-1,1}{
        \draw[\bw, -latex, very thick] (0,0) to (\x,0);
        \draw[\bw, -latex, very thick] (0,0) to (0,\x);
    }
    
    % Gitter
    \foreach \x in {\xmin,...,\xmax}{
        \foreach \y in {\ymin,...,\ymax}{
            \draw[fill] (\x,\y) circle (\r);
            
        }
    }
\end{tikzpicture}
\caption{Konstruktion der ersten Brillouin-Zone}
\label{figure:tikz first bz}
\end{figure}
\begin{figure}[!ht]
\centering
\begin{tikzpicture}[scale=2]   
    \pgfmathsetmacro{\r}{0.03}
    \pgfmathsetmacro{\xmin}{-2}
    \pgfmathsetmacro{\xmax}{2}
    \pgfmathsetmacro{\ymin}{-2}
    \pgfmathsetmacro{\ymax}{2}

    % 2. BZ
    \filldraw[fill=orange!40!white, draw=black] (-1,0) --++ (1,1) --++ (1,-1) --++ (-1,-1) -- cycle;
    
    % 1. BZ
    \filldraw[fill=blue!25!white, draw=black] (-0.5,-0.5) --++ (1,0) --++ (0,1) --++ (-1,0) -- cycle;
    \foreach \x in {-0.5,0.5}{
        \foreach \y in {-0.5,0.5}{
            \draw[\ow, thick] (\x,\y) --+ (-45*\x*\y*4:2cm);
            \draw[\ow, thick] (\x,\y) --+ (135*\x*\y*4:2cm);
        }
    }
    \foreach \v in {-0.5,0.5}{
        \draw[\bw, thick] (\v,\ymin) -- (\v,\ymax);
        \draw[\bw, thick] (\xmin,\v) -- (\xmax,\v);
    }
    \foreach \x in {-1,1}{
        \foreach \y in {-1,1}{
            \draw[\ow, -latex, very thick] (0,0) to (\x,\y);
        }
    }
    
        
    % Gitter
    \foreach \x in {\xmin,...,\xmax}{
        \foreach \y in {\ymin,...,\ymax}{
            \draw[fill] (\x,\y) circle (\r);
            
        }
    }
\end{tikzpicture}
\caption{Konstruktion der zweiten Brillouin-Zone}
\label{figure:tikz second bz}
\end{figure}
\begin{figure}[!ht]
\centering
\begin{tikzpicture}[scale=2]   
    \pgfmathsetmacro{\r}{0.03}
    \pgfmathsetmacro{\xmin}{-2}
    \pgfmathsetmacro{\xmax}{2}
    \pgfmathsetmacro{\ymin}{-2}
    \pgfmathsetmacro{\ymax}{2}
    \def\gw{green!50!white}

    % 3. BZ
    \filldraw[fill=green!25!white, draw=black] (-1,0.5) --++ (0.5,0)  --++ (0,0.5) --++ (1,0) --++ (0,-0.5) --++ (0.5,0) --++ (0,-1) --++ (-0.5,0) --++ (0,-0.5) --++ (-1,0) --++ (0,0.5) --++ (-0.5,0) --++ (0,1) -- cycle;

    % 2. BZ
    \filldraw[fill=orange!40!white, draw=black] (-1,0) --++ (1,1) --++ (1,-1) --++ (-1,-1) -- cycle;
    
    % 1. BZ
    \filldraw[fill=blue!25!white, draw=black] (-0.5,-0.5) --++ (1,0) --++ (0,1) --++ (-1,0) -- cycle;
%    \foreach \x/\y in {1/0.5,1/-0.5,-1/-0.5,-1/0.5}{
%        \draw[\gw, thick] (\x,\y) --+ (90+26.565*\x*\y*2:2cm);
%        \draw[\gw, thick] (\x,\y) --+ (-90+26.565*\x*\y*2:2cm);
%%        \draw[\gw, thick] (\x,\y) --+ (*\x*\y*2:2cm);
%    }
%    \foreach \x/\y in {2/1,1/2,-1/2,-2/1}{
%        \draw[\gw, -latex, very thick] (0,0) to (\x,\y);
%        \draw[\gw, -latex, very thick] (0,0) to (\x,-\y);
%    }
    \foreach \v in {-2,2}{
        \draw[\gw, -latex, very thick] (0,0) to (0,\v);
        \draw[\gw, -latex, very thick] (0,0) to (\v,0);
    }
    \foreach \v in {-0.5,0.5}{
        \draw[\bw, thick] (\v,\ymin) -- (\v,\ymax);
        \draw[\bw, thick] (\xmin,\v) -- (\xmax,\v);
    }
    \foreach \x in {-0.5,0.5}{
        \foreach \y in {-0.5,0.5}{
            \draw[\ow, thick] (\x,\y) --+ (-45*\x*\y*4:2cm);
            \draw[\ow, thick] (\x,\y) --+ (135*\x*\y*4:2cm);
        }
    }
    \foreach \v in {-1,1}{
        \draw[\gw, thick] (\v,0.9*\ymin) -- (\v,0.9*\ymax);
        \draw[\gw, thick] (0.9*\xmin,\v) -- (0.9*\xmax,\v);
    }
    
    
    % Gitter
    \foreach \x in {\xmin,...,\xmax}{
        \foreach \y in {\ymin,...,\ymax}{
            \draw[fill] (\x,\y) circle (\r);
            
        }
    }
\end{tikzpicture}
\caption{Konstruktion der dritten Brillouin-Zone}
\label{figure:tikz third bz}
\end{figure}
\\
Die Zustandsdichte ist gegeben durch
\begin{align}
n(\epsilon) &= \frac{1}{a^2} \sum_{\sigma \textbf{k}} \delta \kla\left( \epsilon - \epsilon_{\sigma}(\textbf{k}) \right) \simeq  \frac{2}{(2\pi)^2} \intr\int_{\mathbb{R}^{2}} \delta \kla\left( \epsilon - \frac{\hbar^2}{2m} \textbf{k}^2 \right) \, \mathrm{d}^{2}k = \\
&=  \frac{2\cdot 2\pi}{(2\pi)^2} \intx\int_{0}^{\infty} \delta\kla\left( \epsilon - u \right) \frac{m}{\hbar^2} \, \mathrm{d}u = \frac{m}{\pi\hbar^2}.
\end{align}
Die Fermi-Energie ist implizit gegeben durch
\begin{align}
\intx\int_{\epsilon_\tecxt\text{min}}^{\epsilon_F} n(\epsilon)\, \mathrm{d}\epsilon &= \frac{N_\tecxt\text{el}}{N} = z_\tecxt\text{el},
\end{align}
wobei $z_\tecxt\text{el}$ die Anzahl an Elektronen pro Elementarzelle ist.
Mit dem gegebenen Wert von $z_\tecxt\text{el} = 2.3\,a^{-2}$ erhalten wir hier
\begin{align}
\epsilon_F &= \frac{\pi\hbar^2z_\tecxt\text{el}}{m} \simeq \SI{5.5e-19}{}a^{-2}.
\end{align}
Um die Dispersion entlang $k_x$ zu skizzieren setzen wir zunächst $k_y=0$ ein
\begin{align}
\epsilon_n(\textbf{k}) &= \frac{\hbar^2}{2m} \klam\left[ \kla\left( k_x - \frac{2\pi}{a}n_x \right)^2 + \frac{4\pi^2}{a^2}n_y^2 \right].
\end{align}
Um die Fermi-Energie in Relation zu setzen, betrachten wir dieses Ergebnis am Rand der ersten Brillouin-Zone.
Dort gilt
\begin{align}
\epsilon_n &= \frac{\hbar^2\pi^2}{2ma^2} \klam\left[ \kla\left( 1 - 2n_x \right)^2 + 4n_y^2 \right] \simeq \SI{3.8e-19}{}a^{-2} \klam\left[ \kla\left( 1 - 2n_x \right)^2 + 4n_y^2 \right].
\end{align}
Im Intervall $k_x \in \klam\left[ 0, \frac{\pi}{a} \right]$ ergibt sich dann für die ersten drei Bänder \autoref{fig:band kx 2d empty}.
\begin{figure}[!ht]
\centering
\includegraphics{band_kx_2d_empty_fermi.pdf}
\caption{Bandstruktur entlang $k_x$ für die ersten drei Bänder im reduzierten Zonenschema.}
\label{fig:band kx 2d empty}
\end{figure} 
Im ausgedehnten Zonenschema ergibt sich einfach die Parabel ausgehend von $n=(0,0)$.
Bei $T=0$ sind alle Zustände bis zur Fermi-Energie besetzt.
Daraus ergibt sich dann die Bedingung
\begin{align}
\epsilon_F &\overset{!}{\ge } \epsilon(\textbf{k}) = \frac{\hbar^2}{2m} \textbf{k}^2,
\end{align}
und somit der Fermi-Impuls 
\begin{align}
k_F &= \sqrt{\frac{2m\epsilon_F}{\hbar^2}} = \sqrt{2\pi z_\tecxt\text{el}} = \frac{2\pi}{a} \sqrt{\frac{a^2z_\tecxt\text{el}}{2\pi}} \simeq \frac{2\pi}{a} \cdot 0.6.
\end{align}
Eine Skizze der Fermi-Fläche für die gegebenen Parameter findet sich in \autoref{fig:tikz fermi surface}.\footnote{Eine schöne Simulation für variable $z_\tecxt\text{el}$ gibt es unter \url{http://lampx.tugraz.at/~hadley/ss2/fermisurface/2d_fermisurface/2dsquare.php}}
\begin{figure}[!ht]
\centering
\begin{tikzpicture}[scale=2]   
    \pgfmathsetmacro{\r}{0.03}
    \pgfmathsetmacro{\zel}{0.605}
    \pgfmathsetmacro{\xmin}{-2}
    \pgfmathsetmacro{\xmax}{2}
    \pgfmathsetmacro{\ymin}{-2}
    \pgfmathsetmacro{\ymax}{2}
    
    % 3. BZ
    \filldraw[fill=green!15!white, draw=black] (-1,0.5) --++ (0.5,0)  --++ (0,0.5) --++ (1,0) --++ (0,-0.5) --++ (0.5,0) --++ (0,-1) --++ (-0.5,0) --++ (0,-0.5) --++ (-1,0) --++ (0,0.5) --++ (-0.5,0) --++ (0,1) -- cycle;

    % 2. BZ
    \filldraw[fill=orange!25!white, draw=black] (-1,0) --++ (1,1) --++ (1,-1) --++ (-1,-1) -- cycle;
    
    % 1. BZ
    \filldraw[fill=blue!15!white, draw=black] (-0.5,-0.5) --++ (1,0) --++ (0,1) --++ (-1,0) -- cycle;
    
    \draw[black, ultra thick, -latex] (0,0) -- (\zel, 0);
    \node[above] at (0.5*\zel, 0) {$k_F$};
    \draw[black, ultra thick] (0,0) circle(\zel);    
    
    % Gitter
    \foreach \x in {\xmin,...,\xmax}{
        \foreach \y in {\ymin,...,\ymax}{
            \draw[fill] (\x,\y) circle (\r);
            
        }
    }
\end{tikzpicture}
\caption{Fermi-Fläche (bzw. Fermi-Ring) für eine Elektronendichte von $z_\tecxt\text{el} = 2.3\,a^{-2}$.}
\label{fig:tikz fermi surface}
\end{figure}
Für ein hinreichend schwaches periodisches Potential kann die Fermi-Energie zumindest in nullter Ordnung noch als konstant angesehen werden.
Insofern führt die vermiedene Kreuzung am Rand der ersten Brillouin-Zone dann dazu, dass ein größerer Teil des ersten Bandes bei $T=0$ besetzt ist. 
In der Skizze würde das dann einer ``Quadratur des Fermi-Kreises'' entsprechen.
\newpage
\section{Eindimensionales Gitter}
Gegen ist ein eindimensionales Gitter mit dem periodischen Potential
\begin{align}
V(x) &= V_0 \cos\kla\left( \frac{2\pi x}{a} \right),
\end{align}
wobei $a$ die Gitterkonstante ist.
In nullter Ordnung ($V\rightarrow 0$) sind die Energien durch
\begin{align}
\epsilon^{(0)} (k) &= \frac{\hbar^2}{2m} \kla\left( k - G \right)^2 \equiv \epsilon_G^{(0)}(k)
\end{align}
gegeben.
Wir berechnen zunächst die Fourier-Koeffizienten des Potentials.
Dabei haben die reziproken Gittervektoren die Struktur $G=\frac{2\pi}{a}n$ für $n\neq n_0\in\mathbb{Z}$.
Es gilt
\begin{align}
V_{G} &= \frac{1}{a}\intx\int_{-a/2}^{a/2} V(x) \e^{-\i Gx}  \, \mathrm{d}x = \frac{V_0}{2a} \intx\int_{-a/2}^{a/2} \kla\left( \e^{\i \kla\left( \frac{2\pi}{a} - G \right)x} + \e^{-\i \kla\left( \frac{2\pi}{a} + G \right)x} \right) \, \mathrm{d}x = \\
&= \frac{V_0}{2a} \intx\int_{-a/2}^{a/2} \kla\left( \e^{\i\frac{2\pi x}{a}(1-n)} + \e^{-\i\frac{2\pi x}{a}(1+n)} \right) \, \mathrm{d}x = \\
&= \frac{V_0}{4\pi} \intx\int_{-\pi}^{\pi} \kla\left( \e^{\i x(1-n)} + \e^{-\i x(1+n)} \right) \, \mathrm{d}x = \frac{V_0}{2} \delta_{|n|,1}
\end{align}
(Im Nachhinein ist das Ergebnis offensichtlich, da die genutzte Darstellung $V(x) \propto \e^{\i \frac{2\pi}{a} x} + \e^{-\i \frac{2\pi}{a} x}$ natürlich bereits der Fouriertransformation entspricht.)
\\
Ohne Bandentartung lautet die Energie in erster Ordnung dann näherungsweise
\begin{align}
\epsilon_{G_0}(k) &\approx \epsilon_{G_0}^{(0)}(k) + \sum_{G\neq G_0} \frac{|V_{G-G_0}|^2}{\epsilon_{G_0}^{(0)}(k) - \epsilon_{G}^{(0)}(k)} = \\
&= \frac{\hbar^2 \kla\left( k - \frac{2\pi}{a}n_0 \right)^2}{2m} + \frac{2m}{\hbar^2} \sum_{n\neq n_0} \frac{\frac{V_0^2}{4} \delta_{|n-n_0|,1}}{\left( k - \frac{2\pi}{a}n_0 \right)^2 - \left( k - \frac{2\pi}{a}n \right)^2} = \\
&= \frac{\hbar^2 \kla\left( k - \frac{2\pi}{a}n_0 \right)^2}{2m} + \frac{2m}{\hbar^2} \frac{V_0^2}{4} \frac{a^2}{(2\pi)^2} \sum_{|\Delta n| = 1} \frac{1}{2\Delta n\left( \frac{ka}{2\pi} - n_0 \right) - (\Delta n)^2} = \\
&= \frac{4\pi^2\hbar^2 \kla\left( \frac{ka}{2\pi} - n_0 \right)^2}{2ma^2} + \frac{2ma^2V_0^2}{16\pi^2\hbar^2} \klam\left[ \frac{1}{2\kla\left( \frac{ka}{2\pi} - n_0 \right) - 1} - \frac{1}{2\kla\left( \frac{ka}{2\pi} - n_0 \right) + 1} \right] = \\
&= \frac{4\pi^2\hbar^2 \kla\left( \frac{ka}{2\pi} - n_0 \right)^2}{2ma^2} + \frac{2ma^2V_0^2}{16\pi^2\hbar^2} \frac{2}{4\kla\left( \frac{ka}{2\pi} - n_0 \right)^2 - 1}.
\end{align}
Da die Summe hier nur zwei Terme enthält, kann die Gleichung auch in der ``besseren'' Näherung noch gelöst werden.
Dabei ergibt sich eine kubische Gleichung, die allerdings ohne etliche Abkürzungen kaum noch aufzuschreiben ist.
Insbesondere ist diese Lösung trotzdem nicht exakt, da wir nur die erste Ordnung der Störungsrechnung betrachten.
%Es gilt
%\begin{align}
%\epsilon_{G_0}(k) &= \epsilon_{G_0}^{(0)}(k) + \sum_{G\neq G_0} \frac{|V_{G-G_0}|^2}{\epsilon_{G_0}(k) - \epsilon_{G}^{(0)}(k)} = \\
%&= \epsilon_{G_0}^{(0)}(k) +  \sum_{n\neq n_0} \frac{\frac{V_0^2}{4} \delta_{|n-n_0|,1}}{\epsilon_{G_0}(k) - \epsilon_{G}^{(0)}(k)} = \\
%&= \epsilon_{G_0}^{(0)}(k) +  \frac{V_0^2}{4} \left[ \frac{1}{\epsilon_{G_0}(k) - \epsilon_{G_0+\Delta G}^{(0)}(k)} + \frac{1}{\epsilon_{G_0}(k) - \epsilon_{G_0 - \Delta G}^{(0)}(k)} \right].
%\end{align}
%Das ist äquivalent zu einer kubischen Gleichung in $\epsilon_{n_0} \equiv \epsilon_{G_0}(k)$ der Form
%\begin{align}
%\kla\left( \epsilon_{n_0} - \epsilon_{n_0}^{(0)} \right) \kla\left( \epsilon_{n_0} - \epsilon_{n_0+1}^{(0)} \right) \kla\left( \epsilon_{n_0} - \epsilon_{n_0-1}^{(0)} \right) &= \frac{V_0^2}{4}\kla\left( \epsilon_{n_0} - \epsilon_{n_0-1}^{(0)} \right) + \frac{V_0^2}{4} \kla\left( \epsilon_{n_0} - \epsilon_{n_0+1}^{(0)} \right).
%\end{align}
%Schreibe $\epsilon := \epsilon_{n_0}$, $e_{k} := \epsilon_{n_0 + k}^{(0)}$ und $\frac{V_0^2}{4} := v$.
%Dann gilt
%\begin{align}
%0 &= \epsilon^3 - \epsilon^2 \kla\left( e_{-1} + e_0 + e_1 \right) + \epsilon \kla\left( e_0e_1 + e_1e_{-1} + e_{-1}e_0 - 2v \right) - e_{-1}e_0e_1 + ve_{-1} + ve_1\equiv\\
%&\equiv \epsilon^3 + b\epsilon^2 + c\epsilon + d.
%\end{align}
%Setze $x := \epsilon + \frac{b}{3}$, $p := c - \frac{b^2}{3}$ und $q := d - \frac{cb}{3} + \frac{2b^3}{27}$. 
%Dann gilt
%\begin{align}
%0 &= x^3 -bx^2 + \frac{b^2x}{3} - \kla\left( \frac{b}{3} \right)^3 + bx^2 - \frac{2b^2x}{3} + \frac{b^3}{9} + cx - \frac{cb}{3} + d = \\
%&= x^3 + px + q.
%\end{align}
%Die Lösung hat die Form
%\begin{align}
%\epsilon &= - \frac{b}{3} + \sqrt[3]{-\frac{q}{2} + \sqrt{\frac{q^2}{4} + \frac{p^3}{27}}} + \sqrt[3]{-\frac{q}{2} - \sqrt{\frac{q^2}{4} + \frac{p^3}{27}}}
%\end{align}

\section{Attraktives Kronig-Penney-Modell}
Gegen ist das eindimensionale periodische Potential
\begin{align}
V(x) &= -v_0\sum_{l\in\mathbb{Z}} \delta(x-la)
\end{align}
mit der Gitterkonstanten $a$.
Das ``atomare'' Problem ist dementsprechend
\begin{align}
E\phi(x) &= -\frac{\hbar^2}{2m}\phi''(x) - v_0\delta(x).
\end{align}
Wir integrieren über ein kleines Intervall $(-\epsilon,\epsilon)$ und finden dann im Grenzwert $\epsilon\rightarrow 0$ die Bedingung
\begin{align}
0 &= -\frac{\hbar^2}{2m}\phi'(x) \Big\vert_{-\epsilon}^\epsilon - v_0,\\
\Rightarrow\ \ \ \phi'(\epsilon) &= \phi'(-\epsilon) - 2\kappa,
\end{align}
wobei $\kappa := \frac{mv_0}{\hbar^2}$.
Für $x\neq 0$ entspricht obiges Problem der freien Schrödingergleichung.
Wir suchen allerdings gebundene Zustände.
Anstatt ebener Wellen setzen wir daher in den beiden Bereichen einen exponentiellen Verlauf ein.
Insbesondere sollte die Lösung symmetrisch in $x$ sein, womit sich der Ansatz auf
\begin{align}
\phi(x) &= A\e^{-k|x|}
\end{align} 
reduziert.
Einsetzen liefert zunächst $E = -\frac{\hbar^2k^2}{2m}$.
Die Ableitung ist $\phi'(x) = -k \frac{x}{|x|}  A\e^{-k|x|}$.
Also muss gelten
\begin{align}
-k\frac{\epsilon}{|\epsilon|} A\e^{-k|\epsilon|} &= -k\frac{-\epsilon}{|-\epsilon|} A\e^{-k|-\epsilon|} - 2\kappa &\Rightarrow && k &= \kappa.
\end{align}
Es existiert also genau ein gebundener Zustand 
\begin{align}
\phi(x) &= \sqrt{\kappa}\e^{-\kappa|x|}
\end{align}
mit Energie $E=-\frac{\hbar^2\kappa^2}{2m} = -\frac{mv_0^2}{2\hbar^2}$.
Insbesondere gibt es also auch nur ein Energieband.
Das Überlappintegral lautet für $R=na\ge 0$
\begin{align}
S(R) &= \int_{\mathbb{R}^{ }} \phi(x-R)\phi(x) \, \mathrm{d}^{ }x = \kappa \int_{\mathbb{R}^{ }} \e^{-\kappa \left\vert x - \frac{na}{2} \right\vert - \kappa \left\vert x + \frac{na}{2} \right\vert} \, \mathrm{d}^{ }x = \\
&= 2\kappa \int_{0}^{na/2} \e^{-\kappa \kla\left( x + \frac{na}{2} + \frac{na}{2} - x \right)}\, \mathrm{d}x + 2\kappa\intx\int_{na/2}^{\infty} \e^{-\kappa \kla\left( x + \frac{na}{2} + x - \frac{na}{2} \right)}\, \mathrm{d}x = \\
&= n\kappa a\e^{-n\kappa a} + 2\kappa \frac{\e^{-2\kappa x}}{-2\kappa}\bigg\vert_{na/2}^\infty = (n\kappa a+1)\e^{-n\kappa a}.
\end{align}
Aufgrund der Symmetrie gilt zudem $S(R) = S(|R|)$.
Wir definieren 
\begin{align}
\Delta V_n(x) &:= -v_0\sum_{l\neq n} \delta(x - la),
\end{align}
was dem Gesamtpotential ohne den Beitrag beim $n$-ten Gitterplatz entspricht.
Das Coulomb-Integral ist dann
\begin{align}
\alpha &= -\intr\int_{\mathbb{R}^{ }} \phi(x)^2\Delta_0(x) \, \mathrm{d}^{ }x = v_0\sum_{l\neq 0} \kappa \intr\int_{\mathbb{R}^{ }} \e^{-2\kappa |x|} \delta(x - la) \, \mathrm{d}^{ }x = \\
&= 2v_0\kappa \sum_{l \ge 1} \e^{-2l\kappa a}.
\end{align}
Die verbleibenden Überlappintegrale sind
\begin{align}
\beta(R) &= -\intr\int_{\mathbb{R}^{ }} \phi(x-R)\phi(x) V_A(x-R) \, \mathrm{d}^{ }x = v_0\kappa\intr\int_{\mathbb{R}^{ }} \e^{-\kappa |x-na|-\kappa |x|} \delta(x - na) \, \mathrm{d}x = \\
&= v_0\kappa \e^{-|n|\kappa a}.
\end{align}
Damit lautet die Gesamtenergie dann
\begin{align}
\epsilon(k) &= E - \frac{\alpha + \sum_{R\neq 0} \e^{-\i kR} \beta(R)}{1 + \sum_{R\neq 0} \e^{-\i kR} S(R)} = \\
&= -\frac{\hbar^2\kappa^2}{2m} - \frac{2v_0\kappa \sum_{n\ge 1}\e^{-2n\kappa a} + \sum_{n\ge 1} 2v_0 \cos(n\kappa a) \kappa \e^{-n\kappa a}}{1 + \sum_{n\ge 1} 2(n\kappa a+1) \cos(n\kappa a) \e^{-n\kappa a}} .%= \\
%&= -\frac{\hbar^2\kappa^2}{2m} - 2v_0\kappa \frac{\frac{1}{\e^{2\kappa a} - 1} +  }{1}
\end{align}
Die geometrischen Reihen können prinzipiell alle berechnet werden, allerdings ist das insbesondere für die Reihe im Nenner recht aufwendig.
%Hier haben wir verschiedene geometrische Reihen verwendet
%\begin{align}
%\sum_{n\ge 1} \cos (nx)\e^{-nx} &= \tecxt\text{Re\,} \sum_{n\ge 1} \e^{(\i-1)nx} = \tecxt\text{Re\,} \frac{1}{\e^{(1-\i) x} - 1} = \\
%&= \tecxt\text{Re\,} \frac{1}{\e^x\cos x - 1 - \i\e^x \sin x} = \frac{\e^x\cos x - 1}{1 + \e^{2x} - 2\e^x \cos(x)}
%\end{align}
In Ordnung $\e^{-\kappa a}$ vereinfacht sich das Ergebnis zu
\begin{align}
\epsilon(k) &\simeq -\frac{\hbar^2\kappa^2}{2m} - 2v_0\kappa \cos(\kappa a) \e^{-\kappa a} \left( 1 - 2[\kappa a+1] \cos(\kappa a) \e^{-\kappa a} \right) \simeq \\
&\simeq -\frac{\hbar^2\kappa^2}{2m} - 2v_0\kappa \cos(\kappa a)\e^{-\kappa a}.
\end{align}

%\begin{align}
%V_{G} &= \frac{1}{a}\intx\int_{-a/2}^{a/2} V(x) \e^{-\i Gx}  \, \mathrm{d}x = 2V_0 \intx\int_{0}^{a/2} \cos\kla\left( \frac{2\pi x}{a} \right) \cos (Gx)\, \mathrm{d}x = \\
%&= 
%\end{align}

\section{Orthorhombisches Gitter}
Gegeben ist das Ergebnis einer vereinfachten Tight-Binding-Näherung für ein einfach orthorhombisches Gitter
\begin{align}
\epsilon (\textbf{k}) &= E_0 - 2\klam\left[ t_a\cos (k_xa) + t_b\cos(k_yb) + t_c\cos(k_zc) \right],
\end{align}
wobei $a$, $b$ und $c$ die Abmessungen der Einheitszelle darstellen.
Der Vektor der Gruppengeschwindigkeit ist definiert als
\begin{align}
\textbf{v}(\textbf{k}) &= \frac{1}{\hbar} \frac{\partial \epsilon}{\partial \textbf{k}} = \frac{2}{\hbar} \begin{pmatrix}
at_a \sin(k_xa) \\ bt_b \sin(k_yb) \\ ct_c \sin(k_zc)
\end{pmatrix}.
\end{align}
Der inverse Tensor der effektiven Masse hat die Komponenten
\begin{align}
\kla\left( \frac{1}{m^\ast} \right)_{jl} &= \frac{1}{\hbar^2} \frac{\partial^2\epsilon}{\partial k_j\partial k_l} = \\
&= \frac{2}{\hbar^2} \frac{\partial}{\partial k_j} \klam\left[ at_a \sin(k_xa) \delta_{lx} + bt_b \sin(k_yb) \delta_{ly} + ct_c \sin(k_zc) \delta_{lz} \right] = \\
&= \frac{2}{\hbar^2} \klam\left[ a^2t_a\cos (k_xa) \delta_{jx}\delta_{lx} + b^2t_b\cos(k_yb) \delta_{jy}\delta_{ly} + c^2t_c\cos(k_zc) \delta_{jz}\delta_{lz} \right].
\end{align}
Der Tensor ist also diagonal.
Explizit lautet er
\begin{align}
\frac{1}{m^\ast} &= \frac{2}{\hbar^2} \begin{pmatrix}
a^2t_a\cos (k_xa) & 0 & 0\\
0& b^2t_b\cos(k_yb) &0 \\
0&0& c^2t_c\cos(k_zc)
\end{pmatrix}.
\end{align}
In der Nähe des $\Gamma$-Punktes gilt $\cos(k_ja_j) \sim 1 - \frac{1}{2} (k_ja_j)^2$, wobei $a_1\equiv a$, $a_2\equiv b$ und $a_3\equiv c$.
Dann vereinfacht sich der Tensor zu
\begin{align}
\kla\left( \frac{1}{m^\ast} \right)_{jl} &\sim \frac{2}{\hbar^2} \delta_{jl} a_j^2t_j\kla\left( 1 - \frac{1}{2}a_j^2k_j^2 \right).
\end{align}
Die effektive Masse hängt hier also sehr direkt mit den drei Raumrichtungen zusammen.
Für $\textbf{k}\rightarrow 0$ gilt insbesondere $m^\ast_{jj} \sim \frac{\hbar^2}{2a_j^2t_j}$.
\\
Im Bandminimum muss gelten $\textbf{v} = 0$.
Das ist hier äquivalent zu $\textbf{k} = 0$, das Minimum entspricht also dem $\Gamma$-Punkt.
Die elektronische Zustandsdichte ist gegeben durch
\begin{align}
n(\epsilon) &= \intr\int_{\mathbb{R}^{3}} \delta\kla\left( \epsilon - \epsilon(\textbf{k}) \right) \, \mathrm{d}^{3}k = \\
&= \frac{1}{abc} \intr\int_{\mathbb{R}^{3}} \delta\kla\left( \epsilon - E_0 + 2t_a\cos(k_x) + 2t_b\cos(k_y) + 2t_c\cos(k_z) \right) \, \mathrm{d}^{3}k \sim\\
&\sim \frac{1}{abc} \intr\int_{\mathbb{R}^{3}} \delta\kla\left( \epsilon - E_0 + 2(t_a+t_b+t_c) - \textbf{k}^2 \right) \, \mathrm{d}^{3}k = \\
&= \frac{4\pi}{abc} \intx\int_{0}^{\infty} \delta\kla\left( \epsilon - E_0 + 2(t_a+t_b+t_c) - u \right) \, u \frac{\tecxt\text{d}k}{\tecxt\text{d}u} \mathrm{d}u = \\
&= \frac{2\pi}{abc} \sqrt{\epsilon - E_0 + 2(t_a+t_b+t_c)}.
\end{align}

\end{document}